%\documentclass[12pt,oneside]{amsart}
\documentclass{scrartcl}

%\documentclass{article}

\usepackage{amsthm,amsmath,amssymb}
\usepackage[numbers]{natbib}

%\usepackage[colorlinks]{hyperref}
\usepackage{hyperref}
\hypersetup{
    colorlinks,%
    citecolor=black,%
    filecolor=black,%
    linkcolor=black,%
    urlcolor=black
}
\usepackage{hypernat}
\usepackage[top=2.5cm, bottom=2.5cm, left=2.8cm,
right=2.8cm]{geometry}
\usepackage{graphicx}


%\setlength{\parskip}{0pt}
%\setlength{\parsep}{0pt}
%\setlength{\headsep}{0pt}
%\setlength{\topskip}{0pt}
%\setlength{\topmargin}{0pt}
%\setlength{\topsep}{0pt}
%\setlength{\partopsep}{0pt}

\linespread{1}

%\usepackage{marvosym}

%\textwidth = 6.5 in \textheight = 9.2 in \oddsidemargin = 0.0 in
%\evensidemargin = 0.0 in \topmargin = -0.0 in \headheight = 0.0 in
%\headsep = 0.0in \parskip = 0.0in \parindent = 0.0in
%\def\baselinestretch{0.871}


\newtheorem{theorem}{Theorem}[section]
\newtheorem{proposition}[theorem]{Proposition}
\newtheorem{problem}[theorem]{Problem}
\newtheorem{fact}[theorem]{Fact}
\newtheorem{lemma}[theorem]{Lemma}
\newtheorem{defn}[theorem]{Definition}
\newtheorem{corollary}[theorem]{Corollary}
\newtheorem{remark}{Remark}[section]
\newtheorem{example}[theorem]{Example}

\newcommand{\E}{{\mathbb E}}
\newcommand{\se}{{\mathcal{E}}}
\newcommand{\R}{{\mathbb R}}
\renewcommand{\P}{{\mathbb P}}
\newcommand{\ac}{{\mathcal{A}}}
\newcommand{\A}{{\mathcal{A}}}
\newcommand{\B}{{\mathcal{B}}}
\newcommand{\C}{{\mathcal{C}}}
\newcommand{\M}{{\mathcal{M}}}
\newcommand{\Mf}{{\mathfrak{M}}}
\newcommand{\T}{{\mathcal{T}}}
\newcommand{\Z}{{\mathcal{Z}}}
\newcommand{\K}{{\mathcal{K}}}
\newcommand{\lc}{{\mathcal{L}}}
\newcommand{\Ps}{{\mathcal{P}}}
\newcommand{\G}{{\mathcal{G}}}
\newcommand{\pa}{{\mathring{p}}}
\newcommand{\F}{{\cal F}}
\newcommand{\samp}{{\mathcal{X}}}
\newcommand{\X}{{\mathcal{X}}}
\newcommand{\hi}{{\hat{\Phi}}}
\newcommand{\data}{{\text{data}}}
\newcommand{\relu}{{\text{ReLU}}}

\newcommand{\ridge}{{\hat{\beta}_{\text{ridge}}(\lambda)}}
\newcommand{\lasso}{{\hat{\beta}_{\text{lasso}}(\lambda)}}
\newcommand{\ridgemu}{{\hat{\mu}_t^{\text{ridge}}(\lambda)}}
\newcommand{\lassomu}{{\hat{\mu}_t^{\text{lasso}}(\lambda)}}


\newcounter{rcnt}[section]
\renewcommand{\thercnt}{(\roman{rcnt})}

\def\qt#1{\qquad\text{#1}}

\def\argmin{\mathop{\rm argmin}}
\def\argmax{\mathop{\rm argmax}}


\setlength{\parskip}{1.4 \medskipamount}
%\sloppy
%\linespread{1.3}

\begin{document}

\title{STAT 238 - Bayesian Statistics  \\
  Lecture Twenty Seven} 
\subtitle{Spring 2026, UC Berkeley}

\author{Aditya Guntuboyina}


\date{06 April 2026}

\maketitle

\section{Gibbs Sampler for Probit Regression}

Our next application of the Gibbs sampler is to Probit
Regression which was introduced in Homework Four. The model is given by
the following. We observe data $(y_1, x_1), \dots, (y_n, x_n)$ with 
  $y_1, \dots, y_n$ binary and $x_1, \dots, x_n \in \R^p$. $x_1,
  \dots, x_n$ are treated non-random and $y_1, \dots, y_n$ random. The 
  parameter is $\beta \in \R^p$. The probit regression model assumes
  that $y_1 \dots, y_n$ are independent given $\beta$ with
  \begin{equation*}
    y_i | \beta \sim \text{Bernoulli} \left(\Phi(x_i' \beta) \right)
  \end{equation*}
  where $\Phi(\cdot)$ is the standard normal cdf. Let $\pi(\cdot)$ be
  our prior for $\beta$. The posterior is then given by:
  \begin{equation*}
    \pi(\beta | \text{data}) \propto \pi(\beta) \prod_{i=1}^n
    \left(\Phi(x_i' \beta) \right)^{y_i} \left(1 - \Phi(x_i' \beta)
    \right)^{1-y_i}. 
  \end{equation*}
  Suppose we take the (improper) uniform prior for $\beta$ so that
  the posterior becomes
  \begin{equation*}
    \pi(\beta | \text{data}) \propto \prod_{i=1}^n
    \left(\Phi(x_i' \beta) \right)^{y_i} \left(1 - \Phi(x_i' \beta)
    \right)^{1-y_i}. 
  \end{equation*}
  It is not actually clear how to use the Gibbs sampler on the above
  posterior. We can try decomposing $\beta$ as $(\beta_1, \dots,
  \beta_p)$ and then attempt to obtain samples from the conditional:
  \begin{equation*}
    \pi(\beta_j | \beta_{-j}, \text{data}) \propto \prod_{i=1}^n
    \left(\Phi(x_i' \beta) \right)^{y_i} \left(1 - \Phi(x_i' \beta)
    \right)^{1-y_i}. 
  \end{equation*}
  I am not sure how to obtain samples from the above conditional.

  There is a trick that people use to run the Gibbs sampler for probit
  regression. Let us introduce random variables $w_1, \dots, w_n$
  which, conditional on $\beta$, are independent with
  \begin{equation*}
   w_i | \beta \sim N(x_i' \beta, 1). 
  \end{equation*}
  We then let
  \begin{equation*}
    y_i = I\{w_i > 0\}. 
  \end{equation*}
  It is then clear that $y_1, \dots, y_n$ defined thus have the same
  likelihood as in the probit regression model. Let $w = (w_1, \dots,
  w_n)$. The posterior of $(w, \beta)$ given the data ($\text{data} =
  y = (y_1, \dots, y_n)$) is
  \begin{align*}
    \pi(w, \beta | \text{data}) &\propto \pi(\beta) \pi(w | \beta)
                                   \pi(y | w, \beta) \\
    &\propto \prod_{i=1}^n \exp \left(\frac{-1}{2} (w_i - x_i'
      \beta)^2\right) \prod_{i=1}^n \left(I\{w_i > 0, y_i = 1\} +
      I\{w_i \leq 0, y_i = 0\} \right) \\
    &\propto \exp \left(\frac{-1}{2} \|w - X \beta\|^2 \right)  \prod_{i=1}^n \left(I\{w_i > 0, y_i = 1\} +
      I\{w_i \leq 0, y_i = 0\} \right). 
  \end{align*}
  where $X$ is, as usual, the design matrix with rows $x_1', \dots,
  x_n'$.

  We shall now run the Gibbs sampler with the decomposition $\theta =
  (w, \beta)$. For this, we need to sample from the conditionals
  $\beta|w, \text{data}$ and $w | \beta, \text{data}$.
  \begin{align*}
    \pi(\beta | w, \text{data}) &\propto \pi(w, \beta | \text{data})
    \\
&\propto \exp \left(\frac{-1}{2} \|w - X \beta\|^2 \right)  \prod_{i=1}^n \left(I\{w_i > 0, y_i = 1\} +
      I\{w_i \leq 0, y_i = 0\} \right) \\
&\propto \exp \left(\frac{-1}{2} \|w - X \beta\|^2 \right) 
  \end{align*}
For a fixed $w \in \R^n$, let 
\begin{equation*}
   \hat{\beta}_w  := \underset{\beta \in \R^p}{\text{argmin}} \|w - X
   \beta\|^2 = (X' X)^{-1} X' w 
\end{equation*}
so that
\begin{equation*}
  \|w - X \beta\|^2 = \|w - X \hat{\beta}_w\|^2 + \|X \hat{\beta}_w -
  X \beta\|^2 = \|w - X \hat{\beta}_w\|^2 + \left(\beta -
    \hat{\beta}_w \right)' (X' X) \left(\beta -
    \hat{\beta}_w \right). 
\end{equation*}
Thus
\begin{equation*}
  \pi(\beta | w, \text{data}) \propto \exp \left(\frac{-1}{2} \left(\beta -
    \hat{\beta}_w \right)' (X' X) \left(\beta -
    \hat{\beta}_w \right) \right). 
\end{equation*}
In other words,
\begin{equation*}
  \beta|w, \text{data} \sim N \left(\hat{\beta}_w, (X' X)^{-1}
  \right). 
\end{equation*}
The next task is to figure out the conditional $w | \beta, \data$:
  \begin{align*}
    \pi(w | \beta, \text{data}) &\propto \pi(w, \beta | \text{data})
    \\
&\propto \exp \left(\frac{-1}{2} \|w - X \beta\|^2 \right)  \prod_{i=1}^n \left(I\{w_i > 0, y_i = 1\} +
         I\{w_i \leq 0, y_i = 0\} \right) \\
& \propto \prod_{i=1}^n \exp \left(\frac{-1}{2} (w_i - x_i' \beta)^2
                                               \right) \left(I\{w_i > 0, y_i = 1\} +
         I\{w_i \leq 0, y_i = 0\} \right)
  \end{align*}
Thus $w_1, \dots, w_n$ are independent conditional on $\beta$ and
$\data$. Further
\begin{equation*}
  \pi(w_i|\beta, \data) = \pi(w_i | \beta, y_i) \propto \exp \left(\frac{-1}{2} (w_i - x_i' \beta)^2
                                               \right) \left(I\{w_i > 0, y_i = 1\} +
         I\{w_i \leq 0, y_i = 0\} \right). 
     \end{equation*}
It makes sense now to separate the two cases $y_i = 1$ and $y_i =
0$. For $y_i = 1$, we get 
     \begin{equation*}
       \pi(w_i | \beta, y_i = 1) \propto \exp \left(\frac{-1}{2} (w_i - x_i' \beta)^2
                                               \right) I\{w_i > 0\}.
                                             \end{equation*}
Thus, conditional on $\beta$ and $y_i = 1$, the random variable $w_i$
has the $N(x_i' \beta, 1)$ distribution trucated to the positive
half-line. 

Here is a fact (from the wikipedia article for truncated normal
distributions:
\url{https://en.wikipedia.org/wiki/Truncated_normal_distribution#Generating_values_from_the_truncated_normal_distribution})
that is useful for simulating from truncated normal 
distributions.
\begin{fact}
  Consider a random variable $X_{\text{trunc}}$  having the truncated
  $N(\mu, \sigma^2)$ truncated to the interval $(a, b)$ (here $a < b$;
  $a$ is allowed to be $-\infty$ and $b$ is allowed to be
  $+\infty$). The density of $X_{\text{trunc}}$ is equal to the
  conditional density of $X$ conditioned on $a < X < b$ where $X \sim
  N(\mu, \sigma^2)$.

  $X_{\text{trunc}}$ can be simulated using the equation:
  \begin{equation*}
    X_{\text{trunc}} = \mu + \sigma \Phi^{-1} \left((1 - U)\Phi(a_0) + U
      \Phi(b_0) \right) ~~\text{where $a_0 := \frac{a - \mu}{\sigma}$
    and $b_0 := \frac{b - \mu}{\sigma}$}
\end{equation*}
and $U$ has the uniform distribution on $(0, 1)$. I will leave as an
exercise for you to verify that the right hand side above indeed has
the required truncated normal distribution. 
\end{fact}
Using this fact, it is straightforward to verify that $w_i|\{\beta,
y_i = 1\}$ can be simulated by
\begin{equation*}
x_i' \beta  + \Phi^{-1}  \left((1 - U) \Phi(-x_i' \beta) + U \right). 
\end{equation*}
This follows directly from the fact above by taking $\mu = x_i'
\beta$, $\sigma = 1$, $a = 0$ and $b = \infty$ (this implies that $a_0
= -x_i' \beta$ and $b_0 = \infty$).

One can argue analogously to deduce that $w_i | \{\beta, y_i = 0\}$
can be simulated by
\begin{equation*}
  x_i' \beta + \Phi^{-1} \left(U \Phi(-x_i' \beta) \right). 
\end{equation*}
The Gibbs sampler for the probit model can be therefore implemented by
the following model: 
\begin{enumerate}
\item Initialize at $\beta^{(0)}$ and $w^{(0)}$.
\item Repeat the following for $t = 1, \dots, N$
  \begin{enumerate}
  \item Pick $J$ to be either $1$ or $2$ with equal probability.
  \item If $J = 1$, take $\beta^{(t+1)} \sim N((X'X)^{-1} X'
    w^{(t)}, (X' X)^{-1})$ and $w^{(t+1)} = w^{(t)}$.
  \item If $J = 2$, take $\beta^{(t+1)} = \beta^{(t)}$ and
    \[
 w_i = 
  \begin{cases} 
x_i' \beta^{(t)}  + \Phi^{-1}  \left((1 - U_i) \Phi(-x_i' \beta^{(t)}) + U_i \right) &
\text{if } y_i = 1 \\
x_i' \beta^{(t)} + \Phi^{-1} \left(U_i \Phi(-x_i' \beta^{(t)}) \right)        &
\text{if } y_i = 0
  \end{cases}
\]
where $U_i \sim Unif(0, 1)$. 
  \end{enumerate}
\end{enumerate}






%\begin{thebibliography}{99}

%\bibitem[Roberts and Rosenthal(2004)]{RobertsRosenthal2004}
%Roberts, G. O. and Rosenthal, J. S. (2004).
%General state space Markov chains and MCMC algorithms.
%\textit{Probability Surveys}, \textbf{1}, 20--71.

%\bibitem[Meyn and Tweedie(2009)]{MeynTweedie2009}
%Meyn, S. P. and Tweedie, R. L. (2009).
%\textit{Markov Chains and Stochastic Stability}, 2nd ed.
%Cambridge University Press.


  
%\bibitem[Chib and Greenberg(1995)]{chib1995}
%Chib, S. and Greenberg, E. (1995).
%Understanding the Metropolis--Hastings algorithm.
%\textit{The American Statistician}, \textbf{49}, 327--335.


%\bibitem[MacKay and Peto(1995)]{MacKay1995HierarchicalDirichlet}
%D.~J.~C. MacKay and L.~C. Peto,
%\newblock ``A hierarchical Dirichlet language model,''
%\newblock \emph{Natural Language Engineering}, vol.~1,
%no.~3, pp.~289--308, 1995.

%\bibitem[Rubin(1981)]{Rubin1981BayesianBootstrap}
%D.~B. Rubin,
%\newblock ``The Bayesian bootstrap,''
%\newblock \emph{The Annals of Statistics}, vol.~9,
%no.~1, pp.~130--134, 1981.
  
%\bibitem[Fisher(1934)]{Fisher1934}
%R.~A. Fisher,
%\newblock ``Probability, likelihood and quantity of information in the logic of
%uncertain inference,''
%\newblock \emph{Proceedings of the Royal Society of London. Series~A,
%Containing Papers of a Mathematical and Physical Character}, vol.~146,
%no.~856, pp.~1--8, 1934.

%\bibitem[Efron(2010)]{Efron}
%Efron, B. (2010)
%\textit{Large-scale inference: empirical Bayes methods for estimation,
%testing, and prediction}. Cambridge University Press, 2010.

%\bibitem[Shumway and Stoffer(2017)]{ShumwayStoffer}
%Shumway, R. H. and Stoffer, D. S. (2017)
%\textit{Time Series Analysis and Its Applications: With R
%  Examples}. Springer, 4th edition. 
  
%     \bibitem[Jaynes(2003)]{Jaynes} Jaynes, E. T, (2003)
%  \textit{Probability theory: the logic of science}. Cambridge
%  University Press, 2003.
 
%  \bibitem[Sinai(1992)]{Sinai} Sinai, Y. G, (1992)
%  \textit{Probability theory: an introductory course}. Springer
%  Verlag, 1992.  

%\bibitem[{deGroot(1986)}]{DeGrootBlackwell}DeGroot, M. H. (1986)
%       A conversation with David Blackwell.
%\newblock \emph{Statistical Science}, \textbf{1}, 40--53.

%         \bibitem[Mosteller(1987)]{Mosteller} Mosteller, F, (1987)
%  \textit{Fifty challenging problems in probability with
%    solutions}. Courier Corporation, 1987.

%    \bibitem[MacKay(2003)]{MacKay} MacKay, D. J. C., (2003)
%  \textit{Information theory, inference, and learning
%  algorithms}. Cambridge University Press, 2003.

%\bibitem[{Lindley(2013)}]{Lindley}Lindley, D. V. (2013)
%   \textit{Understanding Uncertainty}. John Wiley and sons, 2013.


%\end{thebibliography}







\end{document}

%%% Local Variables:
%%% mode: latex
%%% TeX-master: t
%%% End:
